%% ==========================================================================
%% PRISM-Breath -- TRIPOD+AI REPORTING CHECKLIST (supplementary fragment)
%% ==========================================================================
%% Integration note (for the author): \input this file into
%% manuscript/prism_breath_supplementary.tex *before* \end{document}
%% (e.g., right after %% ==========================================================================
%% PRISM-Breath -- EXPANDED APPENDIX ADDITIONS (supplementary material)
%% ==========================================================================
%% Integration note (for the author): these \section and table blocks are meant
%% to be pasted into manuscript/prism_breath_supplementary.tex *before*
%% \end{document}. That file already defines the IEEEtran preamble and the
%% counter renaming (\thetable=S\Roman, \thesection=S\arabic); when pasted in,
%% section and table numbers continue automatically from the existing S1-S4 /
%% Tables SI-SVI. No preamble or \setcounter is needed here. Uses only packages
%% already loaded by the supplementary (amsmath, array, booktabs, amssymb,
%% multirow). Every number below is produced by the cited script; nothing is
%% hand-edited. All results are on the same subject-disjoint five-fold pedrr CV;
%% the 6-model comparisons use the 298-clip set common to all methods (baselines
%% use DeepPhys retrained to 100 epochs, matched compute).
%% ==========================================================================

\section{Full Multi-Metric Comparison}
\label{app:fullmetric}
%% source: scripts/full_metric_comparison.py --deepphys100
Table~\ref{tab:fullmetric} reports all ten evaluation metrics for PRISM-Breath and the five published
extractors on the 298 clips common to every method, with subject-grouped bootstrap 95\% confidence
intervals ($2{,}000$ resamples). Metrics span three families: threshold-free discrimination of tachypnea
(AUROC on the calibrated probability, AUROC ranked by the raw rate--threshold margin, AUROC ranked by the
raw predicted rate, and AUPRC), screening operating points (sensitivity at $90\%$ specificity, specificity
at $90\%$ sensitivity, and the WHO-decision sensitivity/specificity), and regression error (MAE, RMSE).
PRISM-Breath holds the best point estimate on $6$ of $10$ metrics---every discrimination metric plus the
screening-critical specificity-at-$90\%$-sensitivity (where the extractors collapse to $0.05$--$0.21$) and
RMSE. The four metrics the baselines win are sensitivity and Se@Sp$90$ (DeepPhys), raw specificity
(PhysFormer, which achieves it by under-alarming: its sensitivity is only $0.47$), and MAE (DeepPhys,
$5.62$ vs.\ PRISM $6.40$). This dissociation is the paper's thesis in one table: the model with the lowest
regression error is not the best tachypnea screen, because a $\sim$$0.8$~breath/min MAE advantage does not
translate into discrimination near the clinical cutoff.

\begin{table*}[!t]
\caption{Full ten-metric comparison on the common 298-clip five-fold set (subject-bootstrap 95\% CI).
Arrows mark metric direction; \textbf{bold} is the best point estimate in the row. PRISM-Breath is best on
$6/10$ metrics.}
\label{tab:fullmetric}\centering\scriptsize
\resizebox{\textwidth}{!}{%
\begin{tabular}{@{}lcccccc@{}}\toprule
Metric & PRISM-Breath & EfficientPhys & TS-CAN & VIRENet & DeepPhys & PhysFormer\\\midrule
AUROC $\uparrow$          & \textbf{0.874}[0.824,0.920] & 0.755[0.667,0.830] & 0.726[0.642,0.807] & 0.760[0.675,0.837] & 0.777[0.691,0.852] & 0.734[0.662,0.804]\\
AUROC margin $\uparrow$   & \textbf{0.829}[0.765,0.888] & 0.755[0.667,0.830] & 0.726[0.642,0.807] & 0.760[0.675,0.837] & 0.777[0.691,0.852] & 0.734[0.662,0.804]\\
AUROC raw-RR $\uparrow$   & \textbf{0.884}[0.840,0.925] & 0.786[0.718,0.849] & 0.774[0.702,0.838] & 0.795[0.720,0.862] & 0.799[0.727,0.864] & 0.816[0.761,0.867]\\
AUPRC $\uparrow$          & \textbf{0.796}[0.717,0.863] & 0.738[0.648,0.817] & 0.722[0.630,0.803] & 0.756[0.665,0.836] & 0.775[0.689,0.849] & 0.605[0.493,0.721]\\
Se@Sp90 $\uparrow$        & 0.654[0.537,0.797] & 0.704[0.548,0.800] & 0.630[0.530,0.753] & 0.704[0.600,0.805] & \textbf{0.716}[0.620,0.814] & 0.519[0.405,0.628]\\
Sp@Se90 $\uparrow$        & \textbf{0.654}[0.270,0.790] & 0.046[0.000,0.203] & 0.060[0.005,0.159] & 0.051[0.005,0.198] & 0.078[0.004,0.224] & 0.207[0.098,0.345]\\
Sens $\uparrow$           & 0.654[0.548,0.759] & 0.642[0.537,0.744] & 0.630[0.523,0.736] & 0.691[0.588,0.788] & \textbf{0.716}[0.616,0.809] & 0.469[0.355,0.578]\\
Spec $\uparrow$           & 0.908[0.866,0.944] & 0.935[0.898,0.965] & 0.935[0.900,0.967] & 0.940[0.906,0.972] & 0.945[0.910,0.974] & \textbf{0.954}[0.924,0.981]\\
MAE $\downarrow$          & 6.396[5.569,7.369] & 6.102[5.007,7.389] & 6.153[5.002,7.436] & 5.914[4.766,7.219] & \textbf{5.620}[4.588,6.738] & 6.652[5.729,7.654]\\
RMSE $\downarrow$         & \textbf{10.099}[8.412,11.959] & 12.070[9.857,14.435] & 12.314[9.994,14.576] & 11.972[9.628,14.274] & 11.170[9.054,13.267] & 10.996[9.250,12.852]\\\bottomrule
\end{tabular}%
}
\end{table*}

Table~\ref{tab:paired} gives the paired significance tests against PRISM-Breath: the fast-DeLong test on
the AUROC difference and the exact McNemar test on the discordant WHO decisions. PRISM-Breath's AUROC
advantage is significant against all five baselines ($p \le 0.006$). The binary WHO decision is not
significantly different from the extractors (McNemar n.s.), \emph{except against DeepPhys}---whose higher
fixed-threshold sensitivity makes its binary call significantly better ($p{=}0.029$) even though PRISM
out-ranks it on AUROC. This is consistent with the main-paper reading that the advantage lives in
threshold-free discrimination and calibrated operating points rather than the single WHO cutpoint.

\begin{table}[!t]
\caption{Paired tests vs.\ PRISM-Breath (AUROC $0.874$). DeLong on the AUROC gap; exact McNemar on the WHO
tachypnea decision. AUROCs are the canonical values, identical to the multi-metric panel
(Table~\ref{tab:fullmetric}).}
\label{tab:paired}\centering\small
\begin{tabular}{@{}lccc@{}}\toprule
Baseline & AUROC & DeLong $p$ & McNemar $p$\\\midrule
EfficientPhys & 0.755 & $1.2{\times}10^{-3}$ & 0.560 (n.s.)\\
TS-CAN        & 0.726 & $1.9{\times}10^{-4}$ & 0.627 (n.s.)\\
VIRENet       & 0.760 & $1.9{\times}10^{-3}$ & 0.110 (n.s.)\\
DeepPhys      & 0.777 & $6.0{\times}10^{-3}$ & $\mathbf{0.029}$\\
PhysFormer    & 0.734 & $1.3{\times}10^{-4}$ & 0.542 (n.s.)\\\bottomrule
\end{tabular}
\end{table}

\section{Subgroup and Fairness Analysis}
\label{app:subgroup}
%% source: scripts/subgroup_fairness.py --deepphys100
Tables~\ref{tab:agesubgroup}--\ref{tab:sitesubgroup} stratify tachypnea AUROC by clinical age band and by
recording site on the common 298-clip set. Two patterns emerge. First, PRISM-Breath's advantage
concentrates in the \emph{hard} regimes: it leads in the $2$--$12$-month band ($0.871$ vs.\ $0.73$--$0.80$)
and at the two harder, higher-prevalence sites (ICMH, $0.803$ vs.\ $0.64$--$0.71$; Upazila HC, $0.946$ vs.\
$0.79$--$0.87$), while in the easy regimes---the $12$--$60$-month band and the low-prevalence Community
Clinic---all methods cluster near $0.92$--$0.99$ and PRISM's edge is null. Second, the $<2$-month band
remains the hardest: PRISM's discrimination is weak (AUROC $0.684$ on the matched $298$-set, $0.668$ on the
full $332$; $n{=}35$/$38$, wide CI) and its WHO-threshold sensitivity collapses to $0.10$, because the
saturated camera rate ($\sim$$30$) rarely reaches the neonatal $\ge 60$~breath/min cutoff; the CNN
extractors are at or below chance ($0.20$--$0.36$). This band motivates the escalation pathway of
Section~\ref{app:escalation} rather than a screening claim.

\begin{table}[!t]
\caption{Tachypnea AUROC by age band (subject-bootstrap 95\% CI). \textbf{Bold} marks the leader; the
$<2$mo band stays hard (PRISM AUROC $0.684$, WHO-threshold sensitivity $0.10$).}
\label{tab:agesubgroup}\centering\small
\begin{tabular}{@{}lcccc@{}}\toprule
 & overall & 2--12mo & 12--60mo & $<2$mo\\
Method & ($n{=}298$) & ($n{=}127$) & ($n{=}136$) & ($n{=}35$)\\\midrule
PRISM-Breath  & \textbf{0.874} & \textbf{0.871} & 0.939 & \textbf{0.684}\\
EfficientPhys & 0.755 & 0.783 & 0.947 & 0.252\\
TS-CAN        & 0.726 & 0.751 & 0.941 & 0.196\\
VIRENet       & 0.760 & 0.784 & 0.943 & 0.302\\
DeepPhys      & 0.777 & 0.801 & 0.925 & 0.364\\
PhysFormer    & 0.734 & 0.732 & 0.926 & 0.680\\\bottomrule
\end{tabular}
\end{table}

\begin{table}[!t]
\caption{Tachypnea AUROC by recording site (subject-bootstrap 95\% CI). Sites~1--2 are the harder,
higher-prevalence deployments; site~3 is low-prevalence and easy.}
\label{tab:sitesubgroup}\centering\small
\setlength{\tabcolsep}{4pt}
\begin{tabular}{@{}lccc@{}}\toprule
 & ICMH & Upazila HC & Community\\
Method & ($n{=}133$, $32\%$) & ($n{=}78$, $36\%$) & ($n{=}87$, $14\%$)\\\midrule
PRISM-Breath  & \textbf{0.798} & \textbf{0.951} & 0.913\\
EfficientPhys & 0.670 & 0.850 & 0.870\\
TS-CAN        & 0.645 & 0.824 & 0.842\\
VIRENet       & 0.637 & 0.869 & \textbf{0.996}\\
DeepPhys      & 0.706 & 0.831 & 0.923\\
PhysFormer    & 0.690 & 0.794 & 0.852\\\bottomrule
\end{tabular}
\end{table}

\section{Calibration and Selective Prediction}
\label{app:calib}
%% source: scripts/calibration_diagnostics.py
Table~\ref{tab:calib} reports probability-calibration and selective-risk diagnostics on the common
298-clip set. PRISM-Breath's tachypnea probability is roughly twice as well calibrated as the extractors
(ECE $0.070$ vs.\ $0.13$--$0.15$; PhysFormer next at $0.095$), and it has the lowest Brier score ($0.118$)
and the lowest selective-risk AURC ($0.069$ vs.\ $0.16$--$0.20$). Calibration is thus a genuine
differentiator: the reliability-weighted fusion head yields probabilities that support thresholding and
abstention, not just a rank order.

\begin{table}[!t]
\caption{Probability calibration and selective-risk diagnostics (common 298-clip set). ECE/MCE/Brier on the
tachypnea probability; AURC is area under the risk--coverage curve (lower is better throughout).}
\label{tab:calib}\centering\small
\begin{tabular}{@{}lcccc@{}}\toprule
Model & ECE $\downarrow$ & MCE $\downarrow$ & Brier $\downarrow$ & AURC $\downarrow$\\\midrule
PRISM-Breath  & \textbf{0.070} & 0.323 & \textbf{0.118} & \textbf{0.069}\\
EfficientPhys & 0.130 & \textbf{0.284} & 0.131 & 0.187\\
TS-CAN        & 0.145 & 0.295 & 0.138 & 0.201\\
VIRENet       & 0.133 & 0.359 & 0.126 & 0.184\\
DeepPhys      & 0.141 & 0.336 & 0.127 & 0.176\\
PhysFormer    & 0.095 & 0.349 & 0.157 & 0.156\\\bottomrule
\end{tabular}
\end{table}

Table~\ref{tab:conformal} shows that the model-level conformal interval ($\alpha{=}0.20$, target coverage
$0.80$) is well calibrated in aggregate ($0.768$) but age-miscalibrated: it \emph{under}-covers the harder
younger bands ($<2$mo $0.571$; $2$--$12$mo $0.725$) and \emph{over}-covers the easy $12$--$60$-month band
($0.876$) at an essentially constant width ($\sim$$15.7$~breath/min). This motivates the age-conditional
(Mondrian) conformalization used in the escalation pathway (Section~\ref{app:escalation}), where a per-band
$\hat q$ replaces the single width. Per-fold predictive intervals and the hyperparameters are in the
existing supplementary (Tables~SI--SVI); they are not repeated here.

\begin{table}[!t]
\caption{Conformal coverage by age band (nominal $0.80$; single model-level width). ``UNDER''/``OVER'' flag
mis-coverage relative to nominal.}
\label{tab:conformal}\centering\small
\begin{tabular}{@{}lcccl@{}}\toprule
Band & $n$ & Coverage & Width & Status\\\midrule
overall  & 298 & 0.768 & 15.7 & OK\\
2--12mo  & 109 & 0.725 & 15.6 & under\\
12--60mo & 129 & 0.876 & 15.7 & over\\
$<2$mo   & 28  & 0.571 & 15.7 & under\\\bottomrule
\end{tabular}
\end{table}

Table~\ref{tab:riskcov} compares two abstention signals on the risk--coverage curve. Ranking cases by
probability confidence lowers the tachypnea error from $0.161$ at full coverage to $0.061$ when the most
uncertain $40\%$ are deferred (AURC $0.069$); the cross-modal \emph{disagreement} score, by contrast, is a
poor abstention signal (error \emph{rises} to $0.246$; AURC $0.240$). The operational triage that keys on
disagreement therefore auto-clears $62\%$ of cases at accuracy $0.902$ but at a reduced sensitivity of
$0.585$, while the $38\%$ it refers carry sensitivity $0.732$---i.e., disagreement routes many true
positives to referral rather than resolving them, which is why probability-confidence is the preferred
deferral rule.

\begin{table}[!t]
\caption{Risk--coverage: tachypnea error retained at each coverage level under two deferral signals.
Probability-confidence dominates the disagreement score (lower AURC).}
\label{tab:riskcov}\centering\small
\begin{tabular}{@{}lccccc@{}}\toprule
Coverage & 1.0 & 0.9 & 0.8 & 0.7 & 0.6\\\midrule
Disagreement score & 0.161 & 0.179 & 0.197 & 0.225 & 0.246\\
Prob.-confidence    & 0.161 & 0.142 & 0.105 & 0.081 & 0.061\\\midrule
\multicolumn{6}{@{}l}{AURC: disagreement $0.240$ \quad prob.-confidence $\mathbf{0.069}$}\\\bottomrule
\end{tabular}
\end{table}

\section{Decision-Curve Analysis}
\label{app:dca}
%% source: scripts/decision_curve.py --deepphys100
Table~\ref{tab:dca} gives net benefit across risk thresholds $p_t$ for the six models and the treat-all /
treat-none references (common 298-clip set, tachypnea prevalence $0.275$). Screening operates at low $p_t$,
where the cost of a missed infection dominates. In that regime PRISM-Breath has the highest net benefit and
is the only model that exceeds treat-all: at $p_t{=}0.10$ its net benefit is $0.213$ ($95\%$~CI
$[0.16,0.27]$) versus treat-all $0.195$ and the best baseline $0.167$; at $p_t{=}0.20$ it is $0.179$ versus
$0.094$ and $\le 0.151$. As $p_t$ rises past $\sim$$0.30$ (a willingness-to-treat regime less relevant to
screening) some extractors overtake it, mirroring their higher raw sensitivity. PRISM-Breath's net benefit
stays positive and clinically useful across the whole $0.10$--$0.40$ range ($0.213 \to 0.123$).

\begin{table}[!t]
\caption{Decision-curve net benefit vs.\ risk threshold $p_t$ (common 298-clip set, prevalence $0.275$;
subject-bootstrap 95\% CI). \textbf{Bold} marks the best model at each $p_t$.}
\label{tab:dca}\centering\small
\begin{tabular}{@{}lcccc@{}}\toprule
 & \multicolumn{4}{c}{Net benefit at $p_t$}\\\cmidrule(l){2-5}
Strategy & 0.10 & 0.20 & 0.30 & 0.40\\\midrule
Treat-all     & 0.195 & 0.094 & $-0.035$ & $-0.208$\\
Treat-none    & 0.000 & 0.000 & 0.000 & 0.000\\\midrule
PRISM-Breath  & \textbf{0.213} & \textbf{0.179} & \textbf{0.154} & 0.123\\
EfficientPhys & 0.160 & 0.148 & 0.137 & \textbf{0.142}\\
TS-CAN        & 0.152 & 0.135 & 0.128 & 0.122\\
VIRENet       & 0.159 & 0.143 & 0.143 & 0.140\\
DeepPhys      & 0.166 & 0.151 & 0.142 & 0.141\\
PhysFormer    & 0.167 & 0.123 & 0.096 & 0.095\\\bottomrule
\end{tabular}
\end{table}

\section{Threshold-Agnostic Framing}
\label{app:framing}
%% source: scripts/framing_bradypnea.py
The main paper argues that classifying the waveform directly beats regressing a rate and thresholding it.
Two controls here show the effect is not an artifact of the specific WHO tachypnea cutoff. First,
Table~\ref{tab:brady} flips the direction of clinical interest to \emph{bradypnea} ($\text{GT}\le
30$~breath/min, a low-rate operating point matched in prevalence, $0.24$, to the tachypnea base rate): a
ROCKET classifier on each backbone's waveform still beats that backbone's regress-then-threshold framing,
for every backbone, with the gap significant by DeLong ($+0.028$ to $+0.072$ AUROC). The classification
framing therefore helps at both tails, not only for high rates.

\begin{table*}[!t]
\caption{Bradypnea head-swap ($\text{GT}\le 30$~breath/min, prevalence $0.24$). $\Delta$AUROC $=$
classification $-$ regression framing on the \emph{same} backbone; $P(c{>}r)$ is the bootstrap probability
the classifier wins.}
\label{tab:brady}\centering\small
\begin{tabular}{@{}lcccc@{}}\toprule
Backbone & AUROC$_\text{reg}$ & AUROC$_\text{cls}$ & $\Delta$AUROC [95\% CI] & DeLong $p$\\\midrule
EfficientPhys   & 0.925 & 0.953 & $+0.028\,[{+}0.00,{+}0.06]$ & 0.041\\
TS-CAN          & 0.904 & 0.967 & $+0.063\,[{+}0.03,{+}0.10]$ & $1.1{\times}10^{-3}$\\
PRISM (fusion)  & 0.881 & 0.953 & $+0.072\,[{+}0.04,{+}0.11]$ & $3.1{\times}10^{-4}$\\\bottomrule
\end{tabular}
\end{table*}

Second, Table~\ref{tab:sweep} sweeps the decision threshold from $22$ to $64$~breath/min and reports the
framing gain (classification $-$ regression AUROC for the label $\mathbf{1}[\text{GT}>\tau]$) at each. The
gain is positive at every non-degenerate threshold for all three backbones ($\text{frac}(\text{gain}{>}0)=
1.00$). For the extractors it grows monotonically as the threshold moves into the sparse high-rate tail
(EfficientPhys mean $+0.164$, max $+0.339$ at $\tau{=}61$; TS-CAN mean $+0.193$, max $+0.367$ at
$\tau{=}64$)---the reliability-scaling reading: the rarer and higher the positive class, the more the
direct classifier gains over a saturating regressor. For PRISM the gain is smaller and flatter (mean
$+0.097$, largest at the low-rate extreme $\tau{=}22$) because its reliability-weighted fusion head already
recovers most of the discriminative signal, leaving a smaller residual for the head-swap to capture.

\begin{table}[!t]
\caption{Threshold sweep of the framing gain (AUROC$_\text{cls}-$AUROC$_\text{reg}$) for the label
$\mathbf{1}[\text{GT}>\tau]$. ``prev'' is the positive prevalence at $\tau$. Gain is positive at every
threshold for all backbones.}
\label{tab:sweep}\centering\small
\begin{tabular}{@{}ccccc@{}}\toprule
$\tau$ & prev & EfficientPhys & TS-CAN & PRISM (fusion)\\\midrule
22 & 0.96 & $+0.098$ & $+0.060$ & $+0.201$\\
25 & 0.90 & $+0.027$ & $+0.066$ & $+0.126$\\
28 & 0.81 & $+0.040$ & $+0.082$ & $+0.087$\\
31 & 0.73 & $+0.043$ & $+0.073$ & $+0.073$\\
34 & 0.63 & $+0.048$ & $+0.093$ & $+0.023$\\
37 & 0.51 & $+0.084$ & $+0.126$ & $+0.055$\\
40 & 0.42 & $+0.101$ & $+0.143$ & $+0.038$\\
43 & 0.35 & $+0.145$ & $+0.177$ & $+0.055$\\
46 & 0.29 & $+0.176$ & $+0.199$ & $+0.071$\\
49 & 0.25 & $+0.182$ & $+0.218$ & $+0.070$\\
52 & 0.19 & $+0.259$ & $+0.294$ & $+0.110$\\
55 & 0.16 & $+0.284$ & $+0.323$ & $+0.131$\\
58 & 0.13 & $+0.308$ & $+0.328$ & $+0.128$\\
61 & 0.11 & $+0.339$ & $+0.339$ & $+0.143$\\
64 & 0.09 & $+0.324$ & $+0.367$ & $+0.152$\\\midrule
\multicolumn{2}{@{}l}{mean gain} & $+0.164$ & $+0.193$ & $+0.097$\\\bottomrule
\end{tabular}
\end{table}

\section{Neonatal Escalation Pathway}
\label{app:escalation}
%% source: scripts/escalation_pathway.py
Because every camera extractor collapses in the $<2$-month band (Section~\ref{app:subgroup}), the deployed
system does not screen neonates---it \emph{escalates} them. A leakage-free leave-one-fold-out
age-conditional (Mondrian) conformal interval decides, per case, whether the WHO threshold is unambiguously
outside the interval (auto-decide NORMAL/ALERT) or straddles it (REFER to a contact/radar modality).
Table~\ref{tab:escalation} shows the referral rate rises steeply with clinical risk: $79\%$ of $<2$-month
infants are escalated versus $34\%$ of the $12$--$60$-month band. On the cases the system \emph{retains}
(Table~\ref{tab:retained}) it screens at high sensitivity in the older bands ($0.938$ for $12$--$60$mo,
$0.625$ for $2$--$12$mo) and pooled sensitivity $0.735$/specificity $0.993$. Concretely, of the $8$
tachypneic $<2$-month neonates that a naive camera-only screen would clear as normal ($\text{Sens}=0.000$),
escalation routes $6$ to referral, cutting the missed-neonate count from $8$ to $2$.

\begin{table}[!t]
\caption{Escalation (REFER) rate by age band under the leave-one-fold-out age-conditional conformal triage
($n{=}332$). ALERT/NORMAL are the auto-decided cases.}
\label{tab:escalation}\centering\small
\setlength{\tabcolsep}{4pt}
\begin{tabular}{@{}lccccc@{}}\toprule
Band & $n$ & REFER\% & NORMAL\% & ALERT\% & \#tachy\\\midrule
$<2$mo   & 29  & 79.3 & 20.7 & 0.0  & 8\\
2--12mo  & 143 & 48.3 & 44.8 & 7.0  & 49\\
12--60mo & 160 & 34.4 & 55.6 & 10.0 & 29\\
overall  & 332 & 44.3 & 47.9 & 7.8  & 86\\\bottomrule
\end{tabular}
\end{table}

\begin{table}[!t]
\caption{Selective screening on the \emph{retained} (non-escalated) cases. Referral removes the ambiguous
band-straddling cases, leaving a high-specificity screen.}
\label{tab:retained}\centering\small
\begin{tabular}{@{}lccccc@{}}\toprule
Band & $n_\text{ret}$ & pos & Sens & Spec & AUROC\\\midrule
$<2$mo   & 6   & 2  & 0.000 & 1.000 & 0.500\\
2--12mo  & 74  & 16 & 0.625 & 1.000 & 0.871\\
12--60mo & 105 & 16 & 0.938 & 0.989 & 0.985\\
overall  & 185 & 34 & 0.735 & 0.993 & 0.881\\\bottomrule
\end{tabular}
\end{table}

\noindent\textbf{Honest caveats.} The rescue is real but bounded. (i) It requires the \emph{age-conditional}
$\hat q$ (per-band width $\sim$$26$~breath/min for $<2$mo vs.\ $\sim$$6.5$ for $12$--$60$mo); the persisted
single model-level width ($\sim$$8$~breath/min) escalates only $10\%$ of neonates and rescues $0/8$. (ii)
Escalation buys sensitivity with coverage: overall retained coverage is $56\%$ ($44\%$ of cases are
referred rather than screened), which lifts retained sensitivity from the full-coverage $0.663$ to $0.735$.
(iii) Even with age-conditional escalation, $2$ of the $8$ tachypneic neonates are still auto-cleared, so
the pathway mitigates but does not eliminate the neonatal blind spot. A decision-focused end-to-end model
that learns the referral policy directly is the natural next step.
\subsection{Decision-Focused Training Does Not Improve on Calibration Alone}
\label{app:decisionmodel}
%% source: scripts/decision_model_eval.py
The preceding analysis flagged an end-to-end model that learns the referral policy directly as the natural
next step. We built exactly that: PRISM retrained with a differentiable net-benefit (decision-curve)
objective---a soft decision $\sigma(20(p-p_t))$ integrated over the screening threshold grid
$p_t\in[0.05,0.30]$, plus a calibrated abstention gate---under the identical five-fold protocol.
\emph{The decision-focused variant does not improve clinical decision value; if anything it is slightly
worse.} On the common $297$-clip set it lowers AUROC ($0.840$ vs.\ $0.869$) and, by trading sensitivity for
specificity ($0.549/0.949$ vs.\ $0.659/0.907$), reduces net benefit at every operating point in the
screening range (Table~\ref{tab:decmodel}). The paired net-benefit difference at the central threshold
$p_t{=}0.15$ is $-0.019$ ($95\%$ CI $[-0.040,+0.002]$; $P(\text{decision}{>}\text{standard})=0.034$), i.e.\
standard PRISM is favored. The reading is instructive rather than disappointing: once the estimate is
\emph{framed as classification and calibrated} (ECE $0.064$), the model already sits near the
decision-optimal frontier, so an explicit decision-curve loss finds no residual signal and merely shifts the
operating point toward specificity---the wrong direction for a sensitivity-first screen. It is calibration,
not a bespoke decision loss, that delivers the decision-curve advantage.

\begin{table*}[!t]
\caption{Decision-focused training vs.\ standard PRISM (common $297$-clip set, prevalence $0.276$). The
net-benefit objective lowers sensitivity and every net-benefit operating point; standard calibrated PRISM is
favored (paired $\Delta$NB at $p_t{=}0.15$: $-0.019$, $95\%$~CI $[-0.040,+0.002]$).}
\label{tab:decmodel}\centering\small
\begin{tabular}{@{}lccccc@{}}\toprule
 & & & \multicolumn{3}{c}{Net benefit at $p_t$}\\\cmidrule(l){4-6}
Model & AUROC & Sens/Spec & $0.10$ & $0.15$ & $0.20$\\\midrule
Standard PRISM (calibrated) & \textbf{0.869} & $0.659/0.907$ & \textbf{0.214} & \textbf{0.194} & \textbf{0.179}\\
Decision-loss PRISM & $0.840$ & $0.549/0.949$ & $0.206$ & $0.175$ & $0.156$\\\midrule
\multicolumn{6}{@{}l}{\emph{treat-all:} NB$(0.10){=}0.196$, NB$(0.15){=}0.148$, NB$(0.20){=}0.095$; \emph{treat-none:} $0$}\\\bottomrule
\end{tabular}
\end{table*}

\section{Framing Mechanism: Boundary Divergence}
\label{app:mechanism}
%% source: scripts/framing_phase_diagram.py
Why should classifying a waveform and regressing-then-thresholding it ever disagree, given identical inputs?
Table~\ref{tab:mechanism} localizes the effect. On all $332$ subjects the classification framing beats
regression by $+0.051$ AUROC ($0.881$ vs.\ $0.830$; bootstrap $[{+}0.019,{+}0.082]$, $P(\text{gain}{>}0)=
1.00$), but the two framings produce the \emph{same} binary decision on $308/332$ cases. The $24$ decisions
that flip occur almost entirely at the boundary: their mean distance from the cutoff is $2.5$~breath/min
versus $11.9$~breath/min for the agreeing cases. In other words, the framings are interchangeable away from
the threshold and diverge only within $\sim$$2.5$~breath/min of it---exactly where a saturating regressor's
point estimate is least trustworthy and a direct classifier's calibrated margin helps. Stratifying the gain
by estimated distance-to-cutoff confirms it is largest for the mid tertile ($7$--$14$~breath/min,
$+0.117$), i.e.\ the band of genuinely ambiguous cases, and near-zero for cases far from the threshold.

\begin{table*}[!t]
\caption{Boundary-divergence mechanism ($n{=}332$). Top: framing gain by estimated distance from the
cutoff (tertiles of $|\hat r - \tau|$). Bottom: the $24$ decision flips concentrate near the boundary.}
\label{tab:mechanism}\centering\small
\begin{tabular}{@{}lcc@{}}\toprule
Distance-to-cutoff tertile & range (breath/min) & cls$-$reg gain\\\midrule
low  ($n{=}111$) & 0.1--7.0   & $+0.020$\\
mid  ($n{=}110$) & 7.1--14.1  & $+0.117$\\
high ($n{=}111$) & 14.2--37.3 & $+0.089$\\\midrule
\multicolumn{3}{@{}l}{\emph{Decision flips (cls vs.\ reg):} $24/332$ subjects; cls correct $50\%$ of them}\\
\multicolumn{3}{@{}l}{mean $|\hat r - \tau|$: flips $2.5$ vs.\ agreements $11.9$ breath/min}\\\bottomrule
\end{tabular}
\end{table*}

\section{Multiple-comparison control and deployment operating characteristics}\label{sec:sdeploy}
\textbf{Family-wise error control.} The five head-to-head AUROC comparisons against PRISM form a test
family; controlling the false-discovery rate with the Benjamini--Hochberg procedure, all five survive at
$q<0.05$ (Table~\ref{tab:fdr}; largest $q=6.0\times10^{-3}$), so PRISM's superiority is not an artefact of
multiple testing.

\begin{table}[!t]\caption{Benjamini--Hochberg FDR correction across the five DeLong comparisons vs.\ PRISM
($298$ matched subjects). AUROCs as in Table~\ref{tab:paired}.}\label{tab:fdr}\centering\small
\begin{tabular}{@{}lccc@{}}\toprule
Baseline & AUROC & DeLong $p$ & BH $q$\\\midrule
PhysFormer & 0.734 & $1.3\times10^{-4}$ & $4.9\times10^{-4}$\\
TS-CAN & 0.726 & $1.9\times10^{-4}$ & $4.9\times10^{-4}$\\
EfficientPhys & 0.755 & $1.2\times10^{-3}$ & $1.9\times10^{-3}$\\
VIRENet & 0.760 & $1.9\times10^{-3}$ & $2.4\times10^{-3}$\\
DeepPhys & 0.777 & $6.0\times10^{-3}$ & $6.0\times10^{-3}$\\\bottomrule
\end{tabular}\end{table}

\textbf{Deployment PPV/NPV across prevalence.} Predictive values depend on base rate, not the $27.5\%$
cohort prevalence. At a sensitivity-first operating point (Se${\ge}0.90$, Sp${=}0.644$), NPV stays very high
at community-realistic prevalences ($\ge0.98$ for prevalence ${\le}10\%$), so the screen functions as a
strong \emph{rule-out}; PPV is low at low prevalence ($0.12$ at $5\%$)---the unavoidable arithmetic of
screening a rare condition---rising to $0.49$ at cohort prevalence (Table~\ref{tab:ppv}). The
number-needed-to-screen to detect one true case is $22$ at $5\%$ and $11$ at $10\%$ prevalence. A
higher-specificity point (Se${\ge}0.80$, Sp${=}0.782$) raises PPV (e.g.\ $0.29$ at $10\%$) at a sensitivity
cost; the operating point should be set for the target prevalence and the cost of a missed case.

\begin{table}[!t]\caption{Deployment operating characteristics at a sensitivity-first point (Se${\ge}0.90$,
Sp${=}0.644$) across screening prevalence. NNS: subjects screened per true tachypnea case
detected.}\label{tab:ppv}\centering\small
\begin{tabular}{@{}lccc@{}}\toprule
Prevalence & PPV & NPV & NNS\\\midrule
$2\%$ & 0.049 & 0.997 & 55\\
$5\%$ & 0.118 & 0.992 & 22\\
$10\%$ & 0.220 & 0.983 & 11\\
$15\%$ & 0.309 & 0.974 & 7\\
$20\%$ & 0.388 & 0.963 & 6\\
$27.5\%$ (cohort) & 0.490 & 0.946 & 4\\\bottomrule
\end{tabular}\end{table}
). That file already provides
%% the IEEEtran preamble and the S-counter renaming (\thetable=S\Roman,
%% \thesection=S\arabic); section/table numbers here continue automatically.
%% Uses only packages already loaded by the supplementary (amsmath, array,
%% booktabs, amssymb, multirow). No preamble, \usepackage, or \setcounter here.
%% The supplementary is compiled standalone (no bibliography), so TRIPOD+AI is
%% cited by name and the main paper is referenced by section title, not \ref.
%% ==========================================================================

\section{TRIPOD+AI Reporting Checklist}
\label{app:tripod}

The Transparent Reporting of a multivariable prediction model for Individual Prognosis Or Diagnosis --
Artificial Intelligence extension (TRIPOD+AI; Collins \emph{et al.}, \emph{BMJ} 2024) is the reporting
standard for studies that develop or evaluate clinical prediction models built with machine learning.
Table~\ref{tab:tripod-a} through Table~\ref{tab:tripod-e} map every TRIPOD+AI item to the exact location in
this manuscript that addresses it, and record an honest per-item status so that a reader can see at a glance
which requirements are fully met, which are partially met, which are not met, and which require author input
(e.g.\ ethics, funding) that cannot be derived from the manuscript. We deliberately mark items as
\emph{Partial} or \emph{Not met} wherever the evidence is single-cohort, simulation-based, or absent, rather
than overstate the study's clinical validation. Item numbers follow the TRIPOD+AI ordering; authors should
cross-check the final row assignment against the official checklist form before submission.

\noindent\textbf{Status key.} \textbf{Met}~$=$~fully reported in the manuscript;
\textbf{Partial}~$=$~partially addressed, see note; \textbf{Not met}~$=$~a study element not performed or
not reported; \textbf{Author-to-complete}~$=$~administrative content (ethics, funding, registration,
availability statement, accrual dates) not present in the manuscript and to be supplied by the authors before
submission; \textbf{N/A}~$=$~not applicable to this study design.

\begin{table*}[!tp]
\caption{TRIPOD+AI checklist, Part 1 of 5: Title, Abstract, and Introduction.}
\label{tab:tripod-a}\centering\footnotesize
\begin{tabular}{@{}c >{\raggedright\arraybackslash}p{1.65in} >{\raggedright\arraybackslash}p{0.70in} >{\raggedright\arraybackslash}p{3.35in}@{}}
\toprule
Item & TRIPOD+AI item (abbreviated) & Status & Where addressed / notes\\\midrule
\multicolumn{4}{@{}l}{\emph{Title and Abstract}}\\
1  & Title identifies a prediction-model study, target population, and predicted outcome
   & \textbf{Met}
   & Main paper title: non-contact infant tachypnea screening reframed as signal classification (population $=$ infants; outcome $=$ WHO tachypnea).\\
2  & Abstract: structured summary of objectives, data, methods, results
   & \textbf{Met}
   & Main paper Abstract (Goal/Methods/Results/Conclusion): reports AUROC, calibration (ECE), subject-bootstrap CIs, and the neonatal limitation.\\\midrule
\multicolumn{4}{@{}l}{\emph{Introduction}}\\
3a & Background and rationale, including references to existing models
   & \textbf{Met}
   & Main, Sec.~Introduction and Sec.~Related Work (CHROM/POS, DeepPhys, TS-CAN, PhysNet, EfficientPhys, AIR/VIRENet, radar, respiratory audio).\\
3b & Intended use and clinical context; target population
   & \textbf{Met}
   & Main, Sec.~Introduction: frontline WHO/IMCI-style screen runnable from a phone; explicit neonatal-referral policy.\\
4  & Objectives: development, tuning, and/or validation
   & \textbf{Met}
   & Main, Sec.~Introduction (Contributions): develops PRISM-Breath and RespSalience with subject-disjoint internal validation. External validation is a controlled degradation simulation, not a real external video cohort (see Item~5a).\\
\bottomrule
\end{tabular}
\end{table*}

Table~\ref{tab:tripod-b} covers the data, participants, outcome, predictors, sample size, and missing-data
items. The subject-disjoint splitting and leakage handling are strengths of the study; data-provenance
metadata (accrual dates), formal eligibility criteria, and outcome-assessment blinding are the reporting
gaps.

\begin{table*}[!tp]
\caption{TRIPOD+AI checklist, Part 2 of 5: Methods --- data, participants, outcome, predictors, sample size,
missing data.}
\label{tab:tripod-b}\centering\footnotesize
\begin{tabular}{@{}c >{\raggedright\arraybackslash}p{1.65in} >{\raggedright\arraybackslash}p{0.70in} >{\raggedright\arraybackslash}p{3.35in}@{}}
\toprule
Item & TRIPOD+AI item (abbreviated) & Status & Where addressed / notes\\\midrule
5a & Sources of data (development/evaluation), rationale, representativeness
   & \textbf{Partial}
   & Main, Sec.~Datasets $+$ Table ``Datasets'': Pediatric-RR-332 (primary), AIR-125/400, NCHSDB (external), radar phantom $+$ 50 real children. Representativeness limited to a single-site cohort (Discussion). External validation is a camera-realistic \emph{degradation} of NCHSDB polysomnography belts (Sec.~External validation), not a real external video cohort.\\
5b & Dates of participant accrual / data collection
   & \textbf{Author-to-complete}
   & Not stated; authors to add collection dates per cohort.\\
6a & Study setting; number and location of centres
   & \textbf{Partial}
   & Three recording sites analysed (Supp.~Sec.~Subgroup and Fairness); care setting, country, and centre identities not stated.\\
6b & Eligibility criteria for participants
   & \textbf{Partial}
   & Population stated: children aged $<5$~years presenting with cough or difficulty breathing at three
Bangladeshi facilities (Institute of Child and Mother Health, an Upazila Health Complex, a Community
Clinic), stratified by WHO age band ($<$2, 2--12, 12--60~mo) (Main, Sec.~Datasets). The precise
enrolment inclusion/exclusion protocol and any exclusions remain author-to-confirm.\\
6c & Treatments received by participants
   & \textbf{N/A}
   & Screening study; no intervention or treatment.\\
7  & Data preparation, pre-processing, and quality checks
   & \textbf{Met}
   & Main, Sec.~Signal extraction (25~Hz common time base, 0.15--1.5~Hz band-limit); leakage-safe split-then-extract (Supp.~Sec.~Implementation and Reproducibility). Uniformity of pre-processing across skin-tone subgroups not audited (see Item~14).\\
8a & Outcome definition (how and when assessed)
   & \textbf{Met}
   & WHO age-band tachypnea thresholds; ground truth $=$ manual breath counting (Main, Sec.~Introduction/Datasets). \emph{Proxy endpoint:} tachypnea is a screening sign, not diagnosed ARI/pneumonia (Discussion, Limitation~(iii); difficulty-breathing AUROC only $0.59$).\\
8b & Blinding of outcome assessment
   & \textbf{Not met}
   & Blinding of the manual breath-counters to model predictions is not described.\\
9  & Predictors: definition, how and when measured
   & \textbf{Met}
   & Main, Sec.~Signal extraction: monocular-depth respiration proxy, learned optical-flow waveform, audio spectrogram --- all measured automatically per video (no manual predictor assessment to blind).\\
10 & Sample size: how determined and justified
   & \textbf{Partial}
   & Fixed cohort (332 infants); 298-subject matched subset for the paired baseline comparison; subject-disjoint five-fold CV (Main, Sec.~Protocol). No formal power/sample-size calculation; the neonatal subgroup is small ($n{=}29$, wide CIs, stated).\\
11 & Handling of missing data
   & \textbf{Partial}
   & Reliability-/quality-weighted fusion implicitly downweights low-SNR or unavailable modalities (Main, Sec.~PRISM-Breath); no explicit missing-data imputation procedure is reported.\\
\bottomrule
\end{tabular}
\end{table*}

Table~\ref{tab:tripod-c} covers the analytical methods, class imbalance, fairness, model output, and
performance-evaluation items. Partitioning/leakage control, model specification, performance estimation with
uncertainty, and recalibration are fully met; class-imbalance strategy and fairness (specifically skin-tone
equity) are the honest gaps.

\begin{table*}[!tp]
\caption{TRIPOD+AI checklist, Part 3 of 5: Methods --- analytical methods, class imbalance, fairness, model
output, performance evaluation, updating.}
\label{tab:tripod-c}\centering\footnotesize
\begin{tabular}{@{}c >{\raggedright\arraybackslash}p{1.65in} >{\raggedright\arraybackslash}p{0.70in} >{\raggedright\arraybackslash}p{3.35in}@{}}
\toprule
Item & TRIPOD+AI item (abbreviated) & Status & Where addressed / notes\\\midrule
12a & Data partitioning; avoidance of leakage / unwanted resemblance
    & \textbf{Met}
    & Subject-disjoint five-fold, deterministic SHA-256 splits, split-then-extract (Main, Sec.~Protocol; Supp.~Sec.~Implementation). Explicit extractor-leakage audit with a \emph{retracted} TS-CAN result (Main, Sec.~``Extractor robustness and a retracted result'').\\
12b & Handling of predictors in the analysis
    & \textbf{Met}
    & Per-modality encoding (spectral features, temporal scan, self-supervised pretraining) on canonical band-limited signals (Main, Sec.~Method).\\
12c & Model type, all build steps, hyperparameter tuning, internal validation
    & \textbf{Met}
    & PRISM-Breath and RespSalience specified (Main, Sec.~Method); hyperparameters, SSL, and conformal settings (Supp.~Tables~S5--S6); best-validation-checkpoint selection; five-fold internal validation.\\
12d & Handling of heterogeneity in estimates
    & \textbf{Met}
    & Per-fold, per-site, and per-age-band breakdowns (Supp.~Sec.~Per-Fold Results, Subgroup and Fairness).\\
13  & Handling of class imbalance
    & \textbf{Partial}
    & Prevalence reported ($\sim$$0.26$ cohort; $0.155$ NCHSDB) and threshold-free AUROC/AUPRC plus prevalence-dependent PPV/NPV used (Main, Sec.~Calibration/equity; Supp.~Full Multi-Metric Comparison). No explicit resampling/reweighting strategy is stated.\\
14  & Fairness / bias-mitigation approaches
    & \textbf{Partial}
    & Sex (no significant disparity), recording-site, and age-band subgroup analyses reported (Main, Sec.~Calibration/equity; Supp.~Sec.~Subgroup and Fairness). \textbf{Skin-tone/race equity is NOT auditable in a single-population cohort and is NOT mitigated} --- a stated limitation (Discussion, Limitation~(ii)); this specific fairness axis is \emph{Not met}.\\
15  & Specification of the model output
    & \textbf{Met}
    & Outputs $=$ predicted RR, calibrated tachypnea probability (salience mass above the age threshold), and an age-conditional conformal interval (Main, Sec.~RespSalience / conformal triage).\\
16  & Performance measures identified
    & \textbf{Met}
    & RR-MAE/RMSE; AUROC/AUPRC; sensitivity/specificity, Se@Sp90, Sp@Se90; ECE/Brier; decision-curve net benefit (Main, Sec.~Protocol; Supp.~Full Multi-Metric Comparison).\\
17  & Performance estimation with uncertainty
    & \textbf{Met}
    & Subject-level bootstrap 95\% CIs; paired DeLong (AUROC) and exact McNemar (binary decision) (Main, Sec.~Protocol; Supp.~Sec.~Implementation).\\
18  & Model updating / recalibration procedure
    & \textbf{Met}
    & Age-conditional (Mondrian) conformal recalibration replacing a mis-covering pooled interval (Main, Sec.~Age-conditional conformal triage; Supp.~Sec.~Calibration and Selective Prediction, Neonatal Escalation Pathway).\\
\bottomrule
\end{tabular}
\end{table*}

Table~\ref{tab:tripod-d} covers the Results items. Event counts, subgroup performance with confidence
intervals, and updating results are fully reported; a participant-flow diagram, a full demographic table, and
a tabulated deployable model (weights) are the partial items.

\begin{table*}[!tp]
\caption{TRIPOD+AI checklist, Part 4 of 5: Results --- participants, model development, model specification,
performance, updating.}
\label{tab:tripod-d}\centering\footnotesize
\begin{tabular}{@{}c >{\raggedright\arraybackslash}p{1.65in} >{\raggedright\arraybackslash}p{0.70in} >{\raggedright\arraybackslash}p{3.35in}@{}}
\toprule
Item & TRIPOD+AI item (abbreviated) & Status & Where addressed / notes\\\midrule
19a & Flow of participants; numbers with/without outcome (ideally a diagram)
    & \textbf{Partial}
    & Per-fold $n$ and positive counts reported (Supp.~Table~S1: $n{=}332$, 86 positive; NCHSDB $n{=}65$, $1{,}028$ windows). No participant-flow (CONSORT-style) diagram or explicit accounting of exclusions.\\
19b & Participant characteristics (demographics, predictors, outcome)
    & \textbf{Partial}
    & Reported (Main, Sec.~Datasets): $n{=}332$, age $0$--$59$~mo (mean $18$), $195$ male / $137$ female,
three named sites ($137$/$92$/$103$), RR $16$--$96.5$~breath/min, tachypnea prevalence $25.3\%$; per-band
and per-site breakdowns in Supp.~(Subgroup, Equity). A single consolidated demographic table and race/skin
tone are not recorded.\\
19c & Comparison of evaluation-set vs.\ development-set data
    & \textbf{Partial}
    & NCHSDB external cohort described narratively (65 infants, age split, prevalence 15.5\%, non-circular label; Main, Sec.~External validation). Not a side-by-side characteristics comparison.\\
20a & Number of participants and outcome events in each analysis
    & \textbf{Met}
    & 332 full / 298 matched / per-fold counts / NCHSDB windows all reported (Main $+$ Supp.).\\
20b & Unadjusted predictor--outcome associations
    & \textbf{Met}
    & Raw extracted flow (FFT, no encoder) AUROC $0.723$; raw predicted-rate AUROC; full component isolation (Main, Sec.~``Where the performance comes from''; Table ``Component isolation'').\\
21  & Full model presentation enabling individual predictions
    & \textbf{Partial}
    & Architecture and hyperparameters fully specified (Main, Sec.~Method; Supp.~Tables~S5--S6), but trained weights/coefficients are not tabulated; exact reproduction depends on code/checkpoint release (see Item~26b, Author-to-complete).\\
22  & Performance with CIs, including subgroups / fairness
    & \textbf{Met}
    & Tables ``Clinical operating points'', ``PRISM vs.\ extractors'', per-age-band, calibration, decision-curve, and subgroup --- all with 95\% CIs (Main $+$ Supp.).\\
23  & Results of any model updating
    & \textbf{Met}
    & Age-conditional conformal / escalation results: referral rates by band and retained-case screening (Supp.~Sec.~Neonatal Escalation Pathway; Main, Sec.~Age bands).\\
\bottomrule
\end{tabular}
\end{table*}

Table~\ref{tab:tripod-e} covers the Discussion and Other-information items. Interpretation, limitations, and
implications are strengths; the administrative items (registration, funding, ethics/IRB, conflicts, and a
formal data/code-availability statement) are marked Author-to-complete because they are not present in the
manuscript and must not be fabricated.

\begin{table*}[!tp]
\caption{TRIPOD+AI checklist, Part 5 of 5: Discussion and Other information.}
\label{tab:tripod-e}\centering\footnotesize
\begin{tabular}{@{}c >{\raggedright\arraybackslash}p{1.65in} >{\raggedright\arraybackslash}p{0.70in} >{\raggedright\arraybackslash}p{3.35in}@{}}
\toprule
Item & TRIPOD+AI item (abbreviated) & Status & Where addressed / notes\\\midrule
\multicolumn{4}{@{}l}{\emph{Discussion}}\\
24  & Overall interpretation in context of objectives and prior evidence
    & \textbf{Met}
    & Main, Sec.~Discussion: metric reframing, component isolation, and comparison to published $\sim$$2.9$~MAE infant methods.\\
25  & Limitations
    & \textbf{Met}
    & Main, Sec.~Discussion, Limitations (i)--(v): single-site cohort, manual-count ground truth, no external tachypnea \emph{video} cohort, skin-tone equity, proxy endpoint, radar not yet on neonates, retracted leakage result.\\
26a & Intended clinical use and users
    & \textbf{Met}
    & Main, Sec.~Discussion, ``Clinical translation'': a $\ge$2mo triage signal with mandatory neonatal referral and calibrated abstention.\\
26b & Model / code / data accessibility
    & \textbf{Author-to-complete}
    & Reproducibility procedures described (Main, Sec.~Reproducibility; Supp.~Sec.~Implementation; repository scripts). A formal data- and code-availability statement (repository URL, license) is not yet in the manuscript; some data are access-controlled (NCHSDB $=$ PhysioNet credentialed).\\
26c & Barriers to implementation; further research
    & \textbf{Met}
    & Main, Sec.~Discussion: integrated RGB$+$radar device, prospective outcome-anchored multi-site study, real-neonate-in-crib validation.\\\midrule
\multicolumn{4}{@{}l}{\emph{Other information}}\\
27a & Computational environment / software \& hardware (AI item)
    & \textbf{Met}
    & Supp.~Sec.~Implementation and Reproducibility: PyTorch~2.8 (CUDA~12.9), Python~3.13, NVIDIA Quadro RTX~8000; extractor and SSL settings.\\
27b & Study registration and protocol availability
    & \textbf{Author-to-complete}
    & No registration or protocol reference in the manuscript.\\
27c & Funding
    & \textbf{Author-to-complete}
    & No funding statement present.\\
27d & Ethical approval / IRB
    & \textbf{Author-to-complete}
    & No ethics/IRB statement in the manuscript; pediatric human-subjects video requires an approval statement (do not fabricate a number).\\
27e & Conflicts of interest
    & \textbf{Author-to-complete}
    & Not stated.\\
\bottomrule
\end{tabular}
\end{table*}

\noindent\textbf{Validation-status summary and roadmap.} Read together, the checklist shows an internally
validated model with strong leakage control, calibration, and uncertainty reporting, but with clearly bounded
external and clinical validation. The items we mark \emph{Not met} or only \emph{Partial} for scientific
(not administrative) reasons define the remaining work, in priority order: (1)~\textbf{a real external video
cohort with genuine tachypnea} --- present external validation is a camera-realistic \emph{degradation
simulation} on NCHSDB polysomnography belts (Item~5a), not a camera-on-external-cohort result;
(2)~\textbf{prospective / clinical validation} --- the study is retrospective, and the outcome is the WHO
tachypnea proxy rather than a diagnosed ARI/pneumonia endpoint (Items~4, 8a);
(3)~\textbf{a skin-tone / race fairness audit} --- not auditable in the single-population cohort and not yet
mitigated (Item~14); and (4)~\textbf{outcome-assessment blinding} against manual-count reference error
(Item~8b). The \emph{Author-to-complete} items (accrual dates, registration, funding, ethics/IRB, conflicts,
and a formal data/code-availability statement; Items~5b, 26b, 27b--27e) are administrative and must be filled
in by the authors --- with a genuine ethics approval and funding source --- before submission.
